%% ----------------------------------------------------------------
%% Chapter 7 -- Challenges and Solutions
%% ----------------------------------------------------------------
\chapter{Challenges and Solutions}
\label{chap:challenges}

\textbf{Temporal Synchronization.} The two rate (dual rate) nature of both radar (radar sampling frequency) and Wi-Fi (Wi-Fi packet transmission frequency) combined with bursty transmission characteristics create a basic fusion issue. All standard filters rely upon synchronous updates; as such, during FTM gaps, a simple ``hold'' system treats stale (old) Wi-Fi reading values equally to fresh radar data causing estimates to pull backwards and creating the stair-step error. The DCI approach uses velocity extrapolation instead of the prediction of the previous step in order to improve the estimation of the state using the stale sensor information. Additionally, the DCI approach increases the covariance (stale sensor) of the past measurement based on an ``age penalty.'' On the other hand, the neural network accomplishes this by performing a forward fill across the sliding window of data used as input to the network. The neural network identifies values that have been kept constant as fresh using binary flags. Therefore, the neural network is able to learn how quickly the values will degrade.

\textbf{UART Data Integrity under Core Contention.} To ensure precision timing of the Wi-Fi transmission times, the Wi-Fi MAC disables all lower priority interrupts. Since UART byte arrival occurs at a rate of once per 39 microsecond intervals when transmitting at high baud rates, the Wi-Fi interrupt suppression causes the hardware buffer to overflow. This causes corruption of the radar's frame parity structure. To solve this problem, the solution provides for dual-core separation where Core~0 is permanently dedicated to Wi-Fi operation, and Core~1 performs UART data processing and fusion operations. To bridge the two cores, a queue exists which supports zero block sends. As such, if UART data becomes lost during serial ingestion, frames are simply dropped rather than causing the process to stall.

\textbf{Quantization Degradation.} Standard post training quantization rounds off radar range into discrete bins, which results in a geometric noise floor. When integer approximations are made of trigonometric functions and propagated through hidden layers, they can differ substantially from their continuous counterparts. As a result, the regression accuracy of the system is severely compromised. However, because the size of the parameters required to support this architecture is small enough, the entire floating point version of the model fits entirely within internal memory. As such, rapid inference can be performed without loss of performance.

\textbf{Near Field Multipath Distortion.} Below a range of 1.5~m Wi-Fi has significant errors in estimating distance due to saturation effects at the baseband receiver, as well as dense multipath reflections. Standard 20~MHz channels have a high temporal resolution floor which causes the correlator to combine direct and indirect arrivals. Reducing the number of frequency bands in use from 80 to 40 (HT40) halves this ambiguity. In addition, an average-time heuristic applied across 16 burst frames finds the shortest valid path and removes any delayed reflected paths.

\textbf{Radar motion dependency.} The radar processing filter identifies stationary targets as background clutter and terminates tracking on that target. However, a target with respiratory micro-motion will create fractional Doppler shifts during brief pauses; therefore, continuous wireless heartbeats can provide continuity to maintain target identity and coarse proximity until macro-motion occurs again.
